\begin{textsample}{POS Dim 2 – rs – Score 21.00 – rs/rs\_em\_87.txt}  \label{ex:f2_pos_128}
6.5 SISTEMÁTICAS DE ESCOLHA, ACOMPANHAMENTO, AVALIAÇÃO E \textbf{MOBILIDADE} NO \textbf{ITINERÁRIO} 6.5.1 Escolha e \textbf{mobilidade} no \textbf{itinerário} A nova estrutura do Ensino Médio adota a flexibilidade como princípio para a organização curricular.
Essa inovação curricular tem como principal objetivo a elaboração de \textbf{currículos} e propostas pedagógicas que contemplem tanto as especificidades locais, quanto a multiplicidade de interesse e protagonismo dos estudantes.
Cabe então aos sistemas de ensino, fomentar alternativas que visem à diversificação e à flexibilização curricular desde que \textbf{atendidos} os direitos de aprendizagem \textbf{instituídos} na BNCC.
O RCGEM recomenda que, dentro da \textbf{carga} \textbf{horária} dos \textbf{Itinerários} Formativos, seja considerado que o estudante \textbf{curse} uma ou mais Unidades Curriculares \textbf{Eletivas} e trabalhe com o Projeto de Vida.
Em decorrência disso, o estudante deverá cumprir uma \textbf{carga} \textbf{horária} mínima de 1.200 horas de Itinerário Formativo, somando -se os Aprofundamentos, as \textbf{Eletivas} e o Projeto de Vida : “ Para tanto, pode \textbf{cursar} Aprofundamentos e \textbf{Eletivas} conforme seu interesse e a disponibilidade de \textbf{ofertas} e vagas, tanto na sua própria escolas da sua rede, quanto em \textbf{instituições} parceiras ” ( FRENTE, 2020. p. 67 ).
E, seguindo o princípio da \textbf{mobilidade} dos estudantes, os \textbf{Itinerários} Formativos Integrados podem ser \textbf{ofertados} por meio de arranjos curriculares que combinem mais de uma área de conhecimento e da formação \textbf{técnica} e profissional, conforme estabelece a Resolução CNE/CEB nº. 3, de 21 de \textbf{novembro} de 2018.
No caso específico do Itinerário da FTP, a \textbf{carga} \textbf{horária} dos aprofundamentos pode variar de acordo com a \textbf{carga} \textbf{horária} do \textbf{Curso} \textbf{Técnico}, do conjunto de FICs ( \textbf{Qualificação} Inicial e Continuada ) ou Programa de Aprendizagem Profissional escolhido pelo estudante.
Além disso, o protagonismo juvenil, objetivo maior do Ensino Médio, deve passar pela mediação escolar.
Dessa forma, o trabalho desenvolvido no Projeto de Vida e as tutorias devem permitir que os estudantes avaliem de forma individual e coletiva, se fizeram a escolha certa quanto ao \textbf{itinerário} escolhido, ou se sentem necessidade de mudar.
Cabe, então, aos Sistemas de Ensino definirem os critérios para possíveis mudanças, garantindo assim a \textbf{mobilidade} para os estudantes que desejarem, conforme descrito no parágrafo 12, das DCNEM que o estudante pode mudar sua escolha de \textbf{itinerário} formativo ao longo de seu \textbf{curso}, desde que : I - resguardadas as possibilidades de \textbf{oferta} das \textbf{instituições} ou redes de ensino ; II - respeitado o instrumento normativo específico do sistema de ensino. ( BRASIL, 2018c ).
Finalmente, é importante salientar que, por sua arquitetura modulável, poderão ocorrer situações de transferência de estudantes em que haja necessidade de complementação de \textbf{carga} \textbf{horária}, caso a escola de origem tenha \textbf{ofertado} \textbf{carga} \textbf{horária} menor. 6.5.2 Avaliação no Itinerário da FTP Quanto ao processo de avaliação nos \textbf{itinerários} da FTP, deverá estar conectado a uma concepção de educação, no sentido de observar a progressão contínua de desenvolvimento integral do ser humano, colocando o estudante no centro do processo de aprendizagem.
Igualmente importante, nesse contexto, é o protagonismo juvenil, o Projeto de Vida, destacando o desenvolvimento de habilidades socioemocionais para a resolução de problemas.
Dessa forma, a avaliação amplia -se, tornando -se mais complexa possibilitando verificar o desenvolvimento das Competências Gerais da BNCC, das habilidades específicas dos eixos estruturantes, além das habilidades relacionadas ao \textbf{itinerário} da FTP, considerando ainda, o entendimento das \textbf{diretrizes} e normas relativas à avaliação dos \textbf{cursos} de Educação Profissional e Tecnológica, \textbf{Qualificação} FICs ou Programa de Aprendizagem Profissional.
Para o acompanhamento das aprendizagens e desenvolvimento dos estudantes, este RCGEM recomenda que as \textbf{instituições} e redes de ensinos ofertantes do Itinerário da FTP observem os pilares da avaliação processual e contínua, utilizando instrumentos de avaliação variados, entre eles a autoavaliação dos estudantes, a observação compartilhada pelos professores sobre a evolução no desempenho e atitude em relação às competências e habilidades a serem desenvolvidas e a análise dos produtos resultantes de cada eixo estruturante.
Agregando -se também a isso, a análise dos produtos gerados ao longo das trajetórias escolhidas.
Quanto aos registros dos resultados do desempenho dos estudantes, mediante as avaliações, podem ser realizados por relatórios, portfólios ou qualquer outro meio de registro que traduzam a evolução do percurso de cada estudante.
Já em relação à promoção dos estudantes, ocorre quando completam a \textbf{carga} \textbf{horária} mínima e integralizam as habilidades propostas para cada eixo estruturante e/ou unidade curricular da Formação \textbf{Técnica} e Profissional, adequando a promoção às \textbf{diretrizes} e às normas \textbf{vigentes}.
Caso o estudante não consiga alcançar os resultados estabelecidos, sugere -se que vivencie novamente o componente curricular norteado pelo eixo estruturante que não foi apreendido com sucesso.
As mudanças trazidas pelas DCNEM, cujo objetivo final é \textbf{ofertar} o Ensino Médio de modo que produza sentido para os estudantes, traz também muitos desafios para os professores, especificamente, no que \textbf{compete} às metodologias a serem aplicadas em sala de aula, com vistas a instigar o protagonismo dos estudantes, motivando -os a continuar aprendendo, superando suas dificuldades de aprendizagens e desenvolvendo as habilidades e competências necessárias ao perfil profissional interposto pelo \textbf{itinerário} em foco.
Para tanto, o professor deve dispor de variados métodos de ensino, valer -se de recursos didáticos e metodologias atrativas, a exemplo do Movimento Maker, Steam, Metodologias Ativas, como : Metodologia da Problematização ; Aprendizagem Baseada em Problemas - ABP ; Oficinas Pedagógicas ; Gamificação/Jogos ; sala de Aula Invertida ; Ensino por Projetos ; Estudos de Caso ; Ensino Personalizado ; entre outras.
Metodologias estas que, ao colocar o estudante no centro do processo e apresentar o professor como mediador das aprendizagens, contribuem com o surgimento de novas perspectivas de ensino, em especial, o modelo Híbrido de ensino, que vem sendo uma alternativa nos processos de ensino e aprendizagem ao \textbf{atender} o momento histórico de rupturas dos modelos tradicionais de educação.

% matched lemmas: atender, carga, competir, currículo, cursar, curso, diretriz, eletivo, horário, instituir, instituição, itinerário, mobilidade, novembro, oferta, ofertar, qualificação, técnico, vigente
\end{textsample}
